\subsection{Special cases, further results, and open problems}\label{sec:paper-further-results}
This subsection discusses the vanishing of $\SK_2$ beyond Corollary~\ref{cor:second-three-charzero}.
We state conjectures for three-dimensional local fields and for two- and three-dimensional global fields, record a special case, and collect further observations on their scope.
\begin{conjecture}\label{conj:paper-three-local-sk2}
For every three-dimensional local field $F$ and every central division algebra $D$ over $F$, one has $\SK_2(D)=0$.
\end{conjecture}

The conjecture is suggested by the vanishing of $\SK_2(D)$ over local fields of dimension at most two: in dimension zero $D$ is commutative by Wedderburn's theorem, and in dimensions one and two this is Theorems~\ref{thm:paper-local-norm} and~\ref{thm:second-two-vanishing}.
In dimension at least four, the analogous statement fails (Remark~\ref{rem:paper-sk2-dimension-four} below).
In dimension three, Corollary~\ref{cor:second-three-charzero} proves the conjecture when the first residue field has characteristic zero.
The next remark records a mixed-characteristic case with positive first residue characteristic.
The general case is open.

We record several further results, some of which are stated without proof.
We include these statements to make their precise scope available for future work.

\begin{remark}
Let $F$ be a three-dimensional local field with first residue field $k$ of characteristic $p>0$ and with characteristic sequence $(0,p,p)$.
For a uniformizer $\pi$ and a unit $u$ with residue $\bar u\notin k^p$, one has $\SK_2(D)=0$ for every central division algebra $D$ over $F$ of period $p^2$ split by $F(\pi^{1/p},u^{1/p})$.
We omit its proof.
This result is not used elsewhere in the paper.
\end{remark}

\begin{remark}\label{rem:paper-sk2-dimension-four}
The analogous vanishing statement fails in every dimension at least four.
Let $\ell$ be a prime and let $k_0$ be a one-dimensional local field with two linearly disjoint cyclic extensions of degree $\ell$; such fields exist in characteristic zero and in positive characteristic.
Platonov constructed a central division algebra $B$ of index $\ell^2$ over $k_0((x))((y))$ with $\SK_1(B)\simeq\Z/\ell$ \cite[Theorems 4.7 and 5.9]{Platonov1976}.
Over the four-dimensional local field $k_0((x))((y))((t))$, Proposition~\ref{prop:paper-laurent-sk2} shows that $D=B((t))$ has $\SK_2(D)\supseteq\SK_1(B)\cdot\{t\}\ne0$.
For dimension $d>4$, adjoin $d-4$ further central Laurent variables $s$ one at a time: restriction maps $\SK_2$ into $\SK_2$, and the specialization $\partial(\{s\}\cdot-)$ sends the restriction of a class back to that class up to sign, so a nonzero class stays nonzero.
These examples delimit the scope of Conjecture~\ref{conj:paper-three-local-sk2}.
\end{remark}

\begin{remark}\label{rem:paper-wild-cokernel}
The reduced-norm cokernel $C_2(D/F)=\coker(\Nrd_2)$ is cyclic of order $\per[D]$ in dimension three by Theorem~\ref{thm:second-three-coker}.
It is killed by $\deg D$, so it is $p$-primary when $\deg D$ is a power of the characteristic $p$ of the last residue field.
Over iterated Laurent series fields of dimension at least four, this $p$-primary cokernel can be finite or uncountable; we give one example of each, without proof.
\begin{enumerate}
\item Let $F=\Q_2((t_2))((t_3))((t_4))$, a four-dimensional local field with last residue field $\F_2$.
Let $\chi$ be the character of the base-changed extension $F(i)/F$ of the ramified quadratic extension $\Q_2(i)/\Q_2$, and let $D$ be the quaternion division algebra over $F$ with class $(\chi,t_2)$.
The Brauer residue of $[D]$ at $t_2$, computed over $\Q_2((t_2))$ before the outer Laurent variables $t_3,t_4$ are added, is a nonzero character, so $D$ has index two over $\Q_2((t_2))$, and adding Laurent variables preserves this.
Then $C_2(D/F)\simeq(\Z/2)^3$.
\item Let $k=\F_q((u))((v))$, with $q$ a power of a prime $p$, and let $L=k(y)$ with $y^p-v^{p-1}y=u$.
The residue polynomial $X^p-u$ is irreducible, so $L/k$ has ramification index one and purely inseparable residue extension of degree $p$, and the roots $y+cv$ with $c\in\F_p$ make $L/k$ cyclic of degree $p$.
Let $D$ be the base change of the cyclic algebra $(L((t))/k((t)),\sigma,t)$ to the four-dimensional local field $F=k((t))((s))$.
Over $k((t))$, the Brauer residue at $t$ is a nonzero cyclic character, so the cyclic algebra has index $p$.
Adjoining the Laurent variable $s$ preserves the index, so $D$ has index $p$.
Then $C_2(D/F)$ contains $k^\times/N_{L/k}L^\times$ as a direct summand, and this group has the cardinality of $\F_q((u))$; in particular, $C_2(D/F)$ is uncountable.
\end{enumerate}
\end{remark}

\begin{conjecture}\label{conj:paper-global-sk2}
Let $F$ be a two- or three-dimensional global field, that is, a finitely generated field of transcendence degree $h\in\{2,3\}$ over a finite field or of transcendence degree $h-1$ over a number field.
Then $\SK_2(D)=0$ for every central division algebra $D$ over $F$.
\end{conjecture}

\begin{remark}
An $h$-dimensional global field is a finitely generated field of transcendence degree $h$ over a finite field or of transcendence degree $h-1$ over a number field.
We have also proved the following implication for a two-dimensional global field $F$ and a central division algebra $D$ over $F$; we omit the proof:
\[
v_2(\deg D)\le3,\qquad v_3(\deg D)\le2,\qquad
v_\ell(\deg D)\le1\;(\ell\ge5)
\quad\Longrightarrow\quad\SK_2(D)=0.
\]
The cases of indices $8$ and $9$, including characteristics $2$ and $3$, involve Adams operations and Bott lifting.
There are central division algebras with nonzero $\SK_2$ over four-dimensional global fields in both zero and positive characteristic.
Take Platonov's algebra $B$ with $\SK_1(B)\ne0$ over a purely transcendental extension $k_0(x,y)$ of a global field $k_0$ \cite[Theorem 6.2]{Platonov1976}, and choose $0\ne\xi\in\SK_1(B)$.
The projection formula, as in the proof of Proposition~\ref{prop:paper-laurent-sk2}, gives $\xi\cdot\{t\}\in\SK_2(B(t))$, and its image in $\SK_2(B((t)))$ is nonzero by that proposition.
These results are consistent with Conjecture~\ref{conj:paper-global-sk2}, which predicts the vanishing of $\SK_2(D)$ at arbitrary index over two- and three-dimensional global fields.
\end{remark}
